\documentclass[11pt,a4paper]{article}

\usepackage[utf8]{inputenc}
\usepackage[T1]{fontenc}
\usepackage[french]{babel}
\usepackage{lmodern}
\usepackage{microtype}
\usepackage{geometry}
\usepackage{amsmath,amssymb,mathtools}
\usepackage{booktabs}
\usepackage{array}
\usepackage{xcolor}
\usepackage{hyperref}
\usepackage{enumitem}
\usepackage{tcolorbox}

\geometry{margin=2.35cm}
\hypersetup{
  colorlinks=true,
  linkcolor=blue!55!black,
  citecolor=blue!55!black,
  urlcolor=blue!55!black,
  pdftitle={Synthese des travaux confirmes GUP UVIR},
  pdfauthor={Notes de travail}
}
\setlength{\parindent}{0pt}
\setlength{\parskip}{0.6em}
\setlist[itemize]{topsep=0.25em,itemsep=0.2em}
\setlist[enumerate]{topsep=0.25em,itemsep=0.2em}

\newcommand{\dd}{\,\mathrm d}
\newcommand{\lp}{\ell_{\rm P}}
\newcommand{\tP}{t_{\rm P}}
\newcommand{\EP}{E_{\rm P}}
\newcommand{\EH}{E_{\rm H}}
\newcommand{\rhoP}{\rho_{\rm P}}
\newcommand{\rhoDE}{\rho_{\rm DE}}
\newcommand{\rhoC}{\rho_{\rm c}}
\newcommand{\rhoReg}{\rho_{\rm reg}}
\newcommand{\rhoBH}{\rho_{\rm BH}}
\newcommand{\EGUP}{E_{\rm GUP}}
\newcommand{\LH}{L_{\rm H}}
\newcommand{\rhograv}{\rho_{\rm grav}}

\title{Synthese des travaux confirmes autour du cadre GUP--UV/IR\\
\large Resultats mathematiques, contraintes observationnelles et statut physique}
\author{Document de synthese}
\date{22 mai 2026}

\begin{document}
\maketitle

\begin{abstract}
Ce document rassemble les resultats verifies des notes precedentes sur une
cellule de phase deformee de type GUP, sa reformulation position--vitesse, la
regularisation UV de l'energie du vide, les relations Planck--Hubble, et les
tests observationnels publics. Le mot \emph{confirme} est utilise dans un sens
strict : confirme par derivation mathematique dans le modele, par identite
dimensionnelle exacte, ou par contrainte numerique reproductible. Il ne signifie
pas que le GUP soit etabli experimentalement. La conclusion robuste est un
no-go conditionnel : la valeur $\beta_0\sim10^{60}$ qui ajuste directement
l'energie noire ne peut pas etre un parametre universel. La lecture viable est
donc une regularisation UV locale combinee a une projection IR ou
holographique gravitationnelle.
\end{abstract}

\tableofcontents
\newpage

\section{Statut epistemique}

Les resultats sont classes en quatre categories.

\begin{enumerate}
  \item \textbf{Confirme mathematiquement} : formule derivee et verifiee dans
  le cadre de la mesure GUP choisie.
  \item \textbf{Identite exacte} : relation decoulant seulement des definitions
  de $c,\hbar,G,H_0,\Omega_\Lambda$.
  \item \textbf{Contrainte observationnelle conditionnelle} : exclusion ou
  borne obtenue en confrontant le modele a des donnees publiques.
  \item \textbf{Hypothese ouverte} : ingredient necessaire mais non encore
  derive dynamiquement.
\end{enumerate}

\begin{tcolorbox}[colback=yellow!5!white,colframe=yellow!45!black,title=Point cle]
La theorie complete n'est pas confirmee. Ce qui est confirme, ce sont des
consequences conditionnelles : \emph{si} la mesure GUP adoptee est utilisee,
\emph{alors} les formules et exclusions ci-dessous suivent.
\end{tcolorbox}

\section{Mesure GUP confirmee dans le modele}

Le point de depart est la mesure deformee
\begin{equation}
\dd N
=
\frac{g\,\dd^3x\,\dd^3p}
{h^3(1+\beta p^2)^\gamma}.
\end{equation}
Le cas central analyse est
\begin{equation}
\gamma=3,
\qquad
\dd N
=
\frac{g\,\dd^3x\,\dd^3p}
{h^3(1+\beta p^2)^3}.
\label{eq:gup_measure}
\end{equation}

L'integrale maitresse est
\begin{equation}
I_s(\gamma)=
\int_0^\infty
\frac{p^{s-1}\dd p}{(1+\beta p^2)^\gamma}
=
\frac12\beta^{-s/2}
B\left(\frac{s}{2},\gamma-\frac{s}{2}\right),
\end{equation}
valable pour $\gamma>s/2$.

Pour $\gamma=3$, les deux integrales verifiees sont
\begin{align}
\int_0^\infty\frac{p^2\dd p}{(1+\beta p^2)^3}
&=
\frac{\pi}{16\beta^{3/2}},\\
\int_0^\infty\frac{p^3\dd p}{(1+\beta p^2)^3}
&=
\frac{1}{4\beta^2}.
\end{align}

\section{Densite maximale de modes}

La densite totale de modes monoparticulaires est finie :
\begin{equation}
\frac{N_{\max}}{V}
=
\frac{4\pi g}{h^3}
\int_0^\infty
\frac{p^2\dd p}{(1+\beta p^2)^3}.
\end{equation}
Donc
\begin{equation}
\boxed{
n_{\max}
=
\frac{g\pi^2}{4h^3\beta^{3/2}}
=
\frac{g}{32\pi\beta_0^{3/2}\lp^3}.
}
\label{eq:nmax}
\end{equation}

Statut : confirme mathematiquement dans le modele $\gamma=3$.

\section{Energie du vide regularisee}

Pour un champ massless de degenerescence $g$,
\begin{equation}
\rhoReg
=
\frac{4\pi g}{h^3}
\int_0^\infty
\frac{p^2\dd p}{(1+\beta p^2)^3}
\frac12pc.
\end{equation}
Avec l'integrale precedente,
\begin{equation}
\boxed{
\rhoReg
=
\frac{g\pi c}{2h^3\beta^2}
=
\frac{g\rhoP}{16\pi^2\beta_0^2}.
}
\label{eq:rho_reg}
\end{equation}

Statut : confirme mathematiquement comme regularisation UV. Cette formule ne
resout pas seule la constante cosmologique.

\section{Reformulation position--vitesse}

En regime non relativiste, $p=mv$. La mesure devient
\begin{equation}
\dd N
=
\frac{g\,m^3\dd^3x\,\dd^3v}
{h^3(1+\beta m^2v^2)^3}.
\end{equation}
La cellule effective dans l'espace $(x,v)$ est donc
\begin{equation}
\boxed{
\Delta^3x\,\Delta^3v
\sim
\left(\frac{h}{m}\right)^3
(1+\beta m^2v^2)^3.
}
\end{equation}

Le seuil cinematique est
\begin{equation}
\boxed{
v_{\rm GUP}
=
\frac{1}{m\sqrt{\beta}}
=
c\frac{m_{\rm P}}{m\sqrt{\beta_0}}.
}
\label{eq:v_gup}
\end{equation}
La longueur minimale reste
\begin{equation}
\boxed{
\Delta x_{\min}=\hbar\sqrt{\beta}
=
\sqrt{\beta_0}\,\lp.
}
\end{equation}

Statut : confirme comme changement de variables non relativiste. Le seuil
physique exact doit etre formule en impulsion relativiste.

\section{Relation Planck--Hubble--energie noire}

Les definitions sont
\begin{equation}
\LH=\frac{c}{H_0},
\qquad
\rhoC=\frac{3c^2H_0^2}{8\pi G},
\qquad
\rhoDE=\Omega_\Lambda\rhoC.
\end{equation}
Comme
\begin{equation}
\rhoP\lp^2
=
\frac{c^7}{\hbar G^2}\frac{\hbar G}{c^3}
=
\frac{c^4}{G},
\end{equation}
on obtient l'identite exacte
\begin{equation}
\boxed{
\rhoDE
=
\frac{3\Omega_\Lambda}{8\pi}
\rhoP\left(\frac{\lp}{\LH}\right)^2.
}
\label{eq:rho_de_uvir}
\end{equation}

Statut : identite exacte. Elle montre que le facteur $10^{-122}$ est un
rapport d'aires Planck--Hubble, mais elle n'est pas encore une dynamique.

\section{Echelle meV et moyenne geometrique}

L'echelle d'energie associee a l'energie noire est
\begin{equation}
E_{\rm DE}
=
\left[\rhoDE(\hbar c)^3\right]^{1/4}.
\end{equation}
Avec $\EH=\hbar H_0$, on trouve
\begin{equation}
\boxed{
E_{\rm DE}
=
\left(\frac{3\Omega_\Lambda}{8\pi}\right)^{1/4}
\sqrt{\EP\EH}.
}
\label{eq:Ede}
\end{equation}

Numeriquement, avec Planck 2018 et CODATA/NIST 2022 :
\begin{equation}
E_{\rm DE}=2.240365~{\rm meV}.
\end{equation}

\section{Identification directe au vide GUP}

Si l'on impose directement
\begin{equation}
\rhoReg=\rhoDE,
\end{equation}
alors
\begin{equation}
\boxed{
\beta_0^{(\Lambda)}
=
\left(\frac{g}{6\pi\Omega_\Lambda}\right)^{1/2}
\frac{1}{H_0t_P}.
}
\label{eq:beta0_lambda}
\end{equation}

Pour $g=1$ :
\begin{align}
\beta_0^{(\Lambda)}
&=
2.363228\times10^{60},\\
\ell_\Lambda
=
\sqrt{\beta_0^{(\Lambda)}}\,\lp
&=
2.484636\times10^{-5}\ {\rm m},\\
\EGUP^{(\Lambda)}
=
\frac{\EP}{\sqrt{\beta_0^{(\Lambda)}}}
&=
7.941888\ {\rm meV}.
\end{align}

Le rapport avec $E_{\rm DE}$ est exactement
\begin{equation}
\boxed{
\frac{\EGUP^{(\Lambda)}}{E_{\rm DE}}
=
\left(\frac{16\pi^2}{g}\right)^{1/4}
=
\frac{2\sqrt{\pi}}{g^{1/4}}.
}
\label{eq:ratio_egup_ede}
\end{equation}

Statut : confirme mathematiquement et numeriquement. Interpretation physique :
problematique si $\beta_0$ est universel.

\section{Projection UV/IR gravitationnelle}

La borne de non-effondrement d'une region de rayon $L$ donne
\begin{equation}
\boxed{
\rhoBH(L)
=
\frac{3c^4}{8\pi G L^2}
=
\frac{3}{8\pi}\rhoP
\left(\frac{\lp}{L}\right)^2.
}
\label{eq:rho_bh}
\end{equation}
Pour $L=\LH$, $\rhoBH(\LH)=\rhoC$.

Une equation candidate, confirmee algebriquement mais non derivee
dynamiquement, est
\begin{equation}
\boxed{
\rhograv(L)
=
\min\left[
\frac{g\rhoP}{16\pi^2\beta_0^2},
\frac{3c^4}{8\pi G L^2}
\right].
}
\label{eq:min}
\end{equation}

Le croisement des deux branches est
\begin{equation}
\boxed{
L_\times
=
\beta_0\sqrt{\frac{6\pi}{g}}\,\lp.
}
\label{eq:Lcross}
\end{equation}
Pour $\beta_0=1,g=1$ :
\begin{equation}
L_\times=4.341608\,\lp.
\end{equation}

\section{Tests publics : verdicts confirmes}

Le parametre teste est $\beta_0^{(\Lambda)}=2.363228\times10^{60}$ pour
$g=1$. Les tests utilisent des contraintes publiques issues de FIRAS,
Planck+BAO+SNe, SPARC, BBN et etoiles a neutrons.

\begin{table}[h]
\centering
\small
\begin{tabular}{p{0.25\linewidth} p{0.35\linewidth} p{0.28\linewidth}}
\toprule
Test & Resultat numerique & Verdict \\
\midrule
FIRAS/CMB &
$\beta_0\lesssim3.72\times10^{58}$ &
$\beta_0^{(\Lambda)}$ exclu comme photon universel \\
BBN/radiation &
$\beta_0\lesssim2.64\times10^{41}$ &
exclu pour radiation primordiale \\
Etoiles a neutrons &
$\beta_0\lesssim1.66\times10^{38}$ &
exclu pour fermions denses \\
HDE $w$ &
$w(c_{\rm hde}=1)=-0.885$ &
modele testable, $c_{\rm hde}$ doit etre ajuste \\
SPARC acceleration &
$a_{\rm dS}/2\pi=8.63\times10^{-11}{\rm m\,s^{-2}}$ &
accord d'ordre de grandeur seulement \\
\bottomrule
\end{tabular}
\caption{Tests publics du parametre $\beta_0^{(\Lambda)}$.}
\end{table}

Conclusion observationnelle conditionnelle :
\begin{equation}
\boxed{
\beta_0^{(\Lambda)}\sim10^{60}
\text{ ne peut pas etre un parametre universel.}
}
\label{eq:no_go}
\end{equation}

\section{Ce qui est confirme comme no-go}

Les donnees publiques imposent les conclusions suivantes.

\begin{enumerate}
  \item Le GUP $\gamma=3$ regularise bien les integrales UV du vide.
  \item L'identification directe $\rhoReg=\rhoDE$ force un
  $\beta_0\sim10^{60}$.
  \item Cette valeur implique une echelle meV et une longueur effective
  $\sim10^{-5}$ m.
  \item Si cette deformation etait universelle, elle modifierait trop fortement
  les photons du CMB, la radiation primordiale et les fermions denses.
  \item Donc l'energie noire ne peut pas etre simplement la densite brute du
  vide GUP avec un parametre universel.
\end{enumerate}

\section{Ce qui reste ouvert}

Les points suivants ne sont pas confirmes et doivent etre traites comme le coeur
du programme theorique.

\begin{enumerate}
  \item Deriver une mesure GUP covariante relativiste.
  \item Comprendre les signes bosons--fermions et la soustraction du vide plat.
  \item Deriver dynamiquement la projection IR/holographique.
  \item Selectionner la bonne longueur cosmologique $L$.
  \item Obtenir un tenseur effectif $T_{\mu\nu}\simeq-\rho g_{\mu\nu}$.
  \item Produire une prediction precise pour $w(a)$.
\end{enumerate}

\section{Conclusion generale}

La synthese confirmee peut etre condensee ainsi :
\begin{equation}
\boxed{
\text{GUP}
\Rightarrow
\rhoReg<\infty
\quad\text{mais}\quad
\rhoReg\neq\rhoDE
\text{ pour un }\beta_0\text{ universel.}
}
\end{equation}

La direction viable est
\begin{equation}
\boxed{
\text{regularisation UV locale}
\quad+\quad
\text{projection IR/holographique}
\quad\Rightarrow\quad
\rho_{\rm grav}\sim
\rhoP\left(\frac{\lp}{L}\right)^2.
}
\end{equation}

Le resultat le plus fort des travaux n'est donc pas une prediction positive
finale, mais une clarification : le probleme de la constante cosmologique est
un probleme UV/IR. Un cutoff GUP seul ne suffit pas ; une theorie viable doit
expliquer pourquoi seule une partie holographiquement admissible du vide
regularise gravite a grande echelle.

\begin{thebibliography}{9}
\bibitem{CODATA2022}
E. Tiesinga, P. J. Mohr, D. B. Newell and B. N. Taylor,
``CODATA recommended values of the fundamental physical constants: 2022'',
\emph{Reviews of Modern Physics} \textbf{97}, 025002 (2025).

\bibitem{Planck2018}
N. Aghanim et al. (Planck Collaboration),
``Planck 2018 results. VI. Cosmological parameters'',
\emph{Astronomy \& Astrophysics} \textbf{641}, A6 (2020),
arXiv:1807.06209.

\bibitem{FIRAS}
COBE/FIRAS data products, NASA LAMBDA,
\url{https://lambda.gsfc.nasa.gov/product/cobe/firas_products.html}.

\bibitem{PantheonPlus}
Pantheon+SH0ES Data Release,
\url{https://github.com/PantheonPlusSH0ES/DataRelease}.

\bibitem{SPARC}
SPARC galaxy rotation-curve database,
\url{https://zenodo.org/records/16284118}.

\bibitem{CKN}
A. Cohen, D. Kaplan and A. Nelson,
``Effective Field Theory, Black Holes, and the Cosmological Constant'',
\emph{Physical Review Letters} \textbf{82}, 4971 (1999).

\bibitem{LiHDE}
M. Li,
``A Model of Holographic Dark Energy'',
\emph{Physics Letters B} \textbf{603}, 1--5 (2004).

\bibitem{Kempf}
A. Kempf, G. Mangano and R. B. Mann,
``Hilbert space representation of the minimal length uncertainty relation'',
\emph{Physical Review D} \textbf{52}, 1108 (1995).
\end{thebibliography}

\end{document}
